\section{Xác suất và kỳ vọng}
\begingroup\small
\begin{multicols}{2}

\subsection{Xác suất có điều kiện và tính độc lập}
Khi các kết quả trong không gian mẫu $\Omega$ đồng khả năng,
$\Pr(A)=|A|/|\Omega|$; chỉ dùng phép đếm này khi giả thiết
\emph{đồng khả năng} thực sự đúng. Chẳng hạn, chọn đều $k$ phần tử
khác nhau từ $N$ phần tử, trong đó $M$ phần tử được đánh dấu:
\[
\Pr(\text{tất cả đều được đánh dấu})
=\frac{\binom Mk}{\binom Nk}.
\]
Công thức trên chỉ áp dụng khi $\binom Nk>0$.

Với $\Pr(B)>0$,
\[
\Pr(A\mid B)=\frac{\Pr(A\cap B)}{\Pr(B)},\qquad
\Pr(A\cap B)=\Pr(A)\Pr(B\mid A)\quad(\Pr(A)>0).
\]
Nếu $B_1,\ldots,B_m$ là một phân hoạch của $\Omega$ với
$\Pr(B_i)>0$, công thức xác suất toàn phần và Bayes là
\[
\Pr(A)=\sum_i\Pr(A\mid B_i)\Pr(B_i),\quad
\Pr(B_j\mid A)=
\frac{\Pr(A\mid B_j)\Pr(B_j)}{\Pr(A)}
\]
(vế sau yêu cầu $\Pr(A)>0$).
Hai biến cố $A,B$ độc lập khi
$\Pr(A\cap B)=\Pr(A)\Pr(B)$.
\emph{Xung khắc} nghĩa là $A\cap B=\varnothing$; hai khái niệm
này khác nhau. Khi lấy mẫu không hoàn lại, các lần chọn thường
không độc lập.

\subsection{Biến chỉ báo và tính tuyến tính của kỳ vọng}
Với biến ngẫu nhiên rời rạc $X$ có kỳ vọng hữu hạn,
$\mathbb E[X]=\sum_x x\Pr(X=x)$. Đặt $I_A=1$ nếu $A$ xảy ra,
bằng $0$ nếu không; khi đó $\mathbb E[I_A]=\Pr(A)$.
Nếu $X$ đếm số đối tượng thỏa tính chất nào đó, viết
$X=\sum_i I_i$ để được
\[
\mathbb E[X]=\sum_i\Pr(I_i=1).
\]
Tính tuyến tính $\mathbb E[\sum_iX_i]=\sum_i\mathbb E[X_i]$
\emph{không đòi hỏi các biến độc lập}. Ví dụ: chọn đều $k$ trong
$N$ phần tử, $M$ phần tử được đánh dấu. Số phần tử đánh dấu được
chọn có kỳ vọng $M\cdot k/N$.

Nếu $X$ nhận giá trị nguyên không âm và có kỳ vọng hữu hạn,
\[
\mathbb E[X]=\sum_{t\ge1}\Pr(X\ge t).
\]
Công thức này hữu ích khi xác suất của biến cố đạt ít nhất
$t$ dễ tính hơn phân phối của $X$.

\subsection{Quy hoạch động xác suất}
Đặt $f_i(k)$ là xác suất có đúng $k$ lần thành công sau $i$
phép thử Bernoulli \emph{độc lập}, với xác suất thành công ở lần
$i$ là $p_i$. Điều kiện đầu $f_0(0)=1$, còn lại bằng $0$.
Tách theo kết quả lần thử cuối:
\[
f_i(k)=(1-p_i)f_{i-1}(k)+p_i f_{i-1}(k-1).
\]
Duyệt $k$ giảm dần để cập nhật tại chỗ; kết quả có
$\sum_{k=0}^n f_n(k)=1$. Thời gian $O(n^2)$, bộ nhớ $O(n)$.

Tổng quát hơn, nếu trạng thái chuyển theo đồ thị có hướng
không chu trình, $\Pr(v)\mathrel{+}=\Pr(u)p(u,v)$ theo thứ tự
tô pô. Tổng xác suất các cạnh đi ra trạng thái chưa kết thúc phải
bằng $1$. Nếu có chu trình, không được dùng một lượt tô pô.

\subsection{Quy hoạch động và phương trình kỳ vọng}
Gọi $E_s$ là tổng chi phí kỳ vọng từ trạng thái $s$ tới khi kết
thúc; $c_s$ là chi phí ngay tại $s$. Theo bước chuyển đầu tiên,
\[
E_s=c_s+\sum_t p(s,t)E_t,\qquad E_{\rm terminal}=0.
\]
Trên đồ thị không chu trình, tính $E_s$ theo thứ tự tô pô ngược.
Nếu chỉ có cạnh tự quay lại $s$ với xác suất $p(s,s)$, chuyển vế:
\[
E_s=\frac{c_s+\sum_{t\ne s}p(s,t)E_t}{1-p(s,s)}
\quad\text{khi }p(s,s)<1.
\]
Ví dụ mỗi lần thử thành công với xác suất $q>0$, số lần thử kỳ vọng
đến thành công thỏa $E=1+(1-q)E$, nên $E=1/q$.

Với chu trình giữa nhiều trạng thái, lập hệ tuyến tính
$(I-P)E=c$ trên các trạng thái chưa kết thúc. Trước khi giải,
cần xét khả năng không bao giờ tới trạng thái kết thúc: nếu điều
này có xác suất dương và chi phí mỗi bước bị chặn dưới bởi
một hằng số dương, kỳ vọng là
vô hạn. Không chia cho $1-p(s,s)$ nếu vẫn còn phụ thuộc vòng
giữa các trạng thái khác.

\paragraph{Mẫu C++17.}
Hàm thứ nhất tính phân phối số lần thành công; hàm thứ hai tính
kỳ vọng trên đồ thị không chu trình. \texttt{order} là một thứ tự
tô pô hợp lệ, \texttt{cost} là chi phí mỗi trạng thái, và trạng
thái kết thúc có chi phí $0$, không có cạnh đi ra.
Trong \texttt{expectedDAG}, các trạng thái được đánh số $[0,n)$;
\texttt{bernoulliCount} nhận \texttt{p[0]} cho phép thử thứ nhất.
\cppsnippet{code/probability/dp.hpp}

\end{multicols}
\endgroup